% frontmatter/zdania/como-usar.tex -- las instrucciones para el adulto de
% «Zdania», en polaco: las de frontmatter/como-usar.tex, con lo que
% cambia en polaco (las 32 letras de "Pisz po śladzie") y de dónde viene
% la historia (el cuaderno de frases en español, semana a semana).
\section*{Jak korzystać z tego zeszytu}

Ten zeszyt jest dla dziecka, które zna już litery i czyta całe zdania, a
teraz chce czytać coraz płynniej, ćwicząc codziennie po trochu -- nie
uczy czytania od zera. To drugi zeszyt z serii \emph{Uczę się czytać}
po polsku, po \emph{Pierwszych słowach}.

Na każdy dzień od poniedziałku do piątku jest jedna strona, przez cały
rok: także w czasie przerwy świątecznej i wakacji -- właśnie wtedy
najbardziej się przydaje, żeby nie wypaść z rytmu. Każda strona wygląda
zawsze tak samo.

\begin{itemize}[leftmargin=6mm]
\item \textbf{Data} i \textbf{Uwagi}: miejsce, w którym dorosły może
  wpisać datę i, jeśli chce, coś, co zauważył tego dnia -- słowo, które
  sprawiło kłopot, takie, które wyszło świetnie, humor, z jakim dziecko
  usiadło do czytania.
\item \textbf{Zdanie (albo zdania) dnia}, w zielonej ramce. Czyta się je
  na głos, powoli, pokazując palcem każde słowo. Na początku roku jest
  jedno zdanie, na końcu cztery. Tekst wydłuża się z każdą częścią roku,
  a nie co tydzień -- nie ma pośpiechu.
\item Codziennie inne \textbf{zadanie}, zawsze związane z tym, co się
  właśnie przeczytało:
  \begin{itemize}[leftmargin=6mm]
  \item \textbf{\lblDibuja}: narysować coś z tego, o czym mówi zdanie.
  \item \textbf{\lblCompleta}: kilka linii bez kształtu -- resztę
    wymyśla dziecko. Nie ma „dobrego” rysunku.
  \item \textbf{\lblCopia}: przepisać dzisiejsze zdanie na linii. To
    ćwiczenie pisania, a nie wymyślania odpowiedzi -- to przychodzi
    później, kiedy dziecko pisze już samo (od drugiej części jest też
    \textbf{\lblResponde}: krótkie pytanie o to, co się przeczytało,
    żeby sprawdzić, że zostało zrozumiane, a nie tylko przeczytane).
  \item \textbf{\lblRelaciona}: połączyć linią to, co do siebie pasuje.
  \item \textbf{\lblAdivina}: prosta zagadka o historii z tego tygodnia.
  \item \textbf{\lblRodea}, \textbf{\lblVerdaderoFalso} i
    \textbf{\lblBusca}: zakreślić dobrą odpowiedź, zdecydować, czy
    zdanie jest prawdziwe, znaleźć w tekście słowo, o które chodzi.
  \item \textbf{\lblCrea}: zadanie otwarte -- wymyślić, narysować
    własną wersję, opowiedzieć, co by się zrobiło na miejscu bohaterów.
  \item \textbf{\lblRepasa}: co dwa tygodnie, krótkie podsumowanie z
    kratkami do zaznaczenia i zadaniem na koniec.
  \item \textbf{\lblTraza}: mniej więcej co dwa tygodnie, zamiast
    rysunku, prowadzi się palcem albo ołówkiem po kropkach litery --
    wielkiej i małej -- a potem pisze się ją samodzielnie na linii. Przez
    cały rok wszystkie 32 litery polskiego alfabetu, po kolei, od
    \emph{a} do \emph{ż}.
  \end{itemize}
\end{itemize}

Strony opowiadają historię \textbf{Lucíi}, jej brata \textbf{Daniego},
ich psa \textbf{Toby'ego} i reszty rodziny. Bohaterowie są z nami przez
cały rok, więc każdy tydzień to nowy rozdział tej samej historii -- tej
samej, tydzień po tygodniu, co w hiszpańskim zeszycie ze zdaniami: jeśli
w domu jest rodzeństwo, każde może czytać swój zeszyt, o tym samym.
Rodzina mieszka w Hiszpanii, więc w historii są też hiszpańskie
zwyczaje: pieczone kasztany pierwszego listopada, Trzech Króli,
karnawał.

Nie trzeba kończyć zadania tego samego dnia, jeśli nie ma czasu --
najważniejsze jest czytanie. Wystarczy pięć albo dziesięć minut
dziennie.

Ten plan jest pomyślany dla „średniego” tempa. Jeśli przez cały tydzień
dziecko czyta stronę bez żadnego błędu za pierwszym razem, nie trzeba
czekać do następnej części roku -- można czytać więcej niż jedną stronę
dziennie. Jeśli zacina się prawie na każdym słowie, niech przeczyta tylko
pierwsze zdanie, a resztę przeczyta na głos dorosły: najważniejsze jest,
żeby czytanie było przyjemnością, a nie walką.

Na koniec każdej części jest mały medal, nie tylko dyplom na ostatni
dzień -- cel bliższy niż 260 dni. Na końcu zeszytu są też
\textbf{odpowiedzi} do zadań \lblAdivina\ i \lblRelaciona.
\newpage
